\documentclass{article}
\usepackage{amsmath,amssymb}
\usepackage[vocab={queues}]{leanify}
\newtheorem{lemma}{Lemma}
\begin{document}

\section{Notation}

We write $\rho(\lambda,\mu)$ for the traffic intensity of the queue with arrival rate $\lambda$ and service rate $\mu$.
We write $L(\lambda,\mu)$ for the mean number of customers in the M/M/1 queue with arrival rate $\lambda$ and service rate $\mu$.
We write $L_q(\lambda,\mu)$ for the mean number of customers waiting in the M/M/1 queue with arrival rate $\lambda$ and service rate $\mu$.
We write $W(\lambda,\mu)$ for the mean sojourn time in the M/M/1 queue with arrival rate $\lambda$ and service rate $\mu$.
We write $P_n(\lambda,\mu)$ for the probability of $n$ customers in the M/M/1 queue with arrival rate $\lambda$ and service rate $\mu$.
We write $B(n,a)$ for the Erlang B formula for $n$ servers and offered load $a$.
We write $C(n,a)$ for the Erlang C formula for $n$ servers and offered load $a$.

\section{The M/M/1 queue}

\begin{lemma}\label{example}
$\rho(1,2)=\frac12$, $L(1,2)=1$, $L_q(1,2)=\frac12$, $W(1,2)=1$ and $P_0(1,2)=\frac12$.
\end{lemma}

\begin{lemma}\label{little}
Let $\lambda,\mu$ be real numbers with $0<\lambda<\mu$. Then $L(\lambda,\mu)=\lambda W(\lambda,\mu)$.
\end{lemma}

\section{Erlang's formulas}

\begin{lemma}\label{erlang}
$B(1,1)=\frac12$ and $C(2,1)=\frac13$.
\end{lemma}

\end{document}
